\lecture{3}{Tue 16 Jan 2024 10:30}{Last day linear algebra!}



\subsection{Technicalities on bra-ket notation}

\subsubsection{Kets and Bras}



The \textbf{bra-ket} notation represents inner product: $\left\langle x \middle| y \right\rangle $. The  \textbf{kets} are just notations for vectors in $\mathscr{H}$.
For example, if $A: \mathscr{H} \to  \mathscr{H}$ is a Hermitian operator with an eigenvalue $\lambda$, we write $\left| \lambda \right\rangle $ for an eigenvector of $A$ with eigenvalue $\lambda$. Thus,  $A \left| \lambda \right\rangle  = \lambda \left| \lambda \right\rangle $.

The bra-ket notation can also be understood as $\left\langle x \middle| y \right\rangle := \left\langle \left| x \right\rangle , \left| y \right\rangle  \right\rangle $.

The \textbf{bras} are elements in the dual space: $\mathscr{H}^*$, identifiable with $\mathscr{H}$. If $\mathscr{H}$ is a Hilbert space, then the \textit{dual Hilbert space} $\mathscr{H}^*$ is defined as
\[
	\mathscr{H}^* = \{\rho: \mathscr{H} \to  \mathbb{C} | \rho \text{ linear and bounded}\} 
.\] 
\begin{remark}
	For infinite-dimensional Hilbert spaces $\mathscr{H}$ the property of $\rho$ being a \textit{bounded} linear functional becomes important.
\end{remark}

It is easy to check that $\mathscr{H}^*$ is a vector space. To make it a Hilbert space, we need to introduce an inner product on $\mathscr{H}^*$. 

Given $\rho \in \mathscr{H}^*$, let us define the \textit{length} (or \textit{operator norm}) of $\rho$ to be
\[
	\|\rho\| := \inf \{\|\rho \left| \psi \right\rangle \| / \|\left| \psi \right\rangle \|: \left| \psi \right\rangle \neq 0, \left| \psi  \right\rangle \in \mathscr{H}\} 
.\] 
\begin{remark}
	This norm endows $\mathscr{H}^*$ with the structure of a \textit{Banach space}. Banach spaces are similar to Hilbert spaces in that we have a good notion of length in both of them; but Banach spaces does not have a proper notion of angle.
\end{remark}


In fact, $\|\cdot \|_{\mathscr{H}^*}$ satisfies the \textit{polarization identity} which makes it unique.

\subsubsection{Canonical identification with the dual}

There is a \textit{canonical map} between $\mathscr{H}$ and $\mathscr{H}^*$ :
\begin{align*}
	\Phi: \mathscr{H} &\longrightarrow \mathscr{H}^* \\
	\left| \psi \right\rangle  &\longmapsto \Phi(\left| \psi \right\rangle ) = \left\langle \psi \right| 
\end{align*}
where
\begin{align*}
	\left| \psi \right\rangle : \mathscr{H} &\longrightarrow \mathbb{C} \\
	\left| \phi \right\rangle  &\longmapsto \left\langle \psi \middle| \phi \right\rangle 
.\end{align*}

On HW1, we have proved that  $\Phi$ is a bijection. Note that $\Phi$ is \textbf{not} linear; it is actually \textbf{conjugate linear}: 
\[
	\Phi (c \left| \psi \right\rangle ) \left| \phi \right\rangle = c^* \left\langle \psi \middle| \phi \right\rangle 
.\] 

\begin{remark}
	\textit{Projective symmetries} matters in quantum mechanics, this is tightly connected with the fact that $\mathscr{H} \cong \mathscr{H}^*$. More on that in future.
\end{remark}

\subsubsection{Operators in Bra-Ket Notation}

The bra-ket notation is especially useful because we can combine bras and kets in the ``wrong order'' and get something helpful. In the ``correct order'' we get the \textit{inner product}
\[
	\mathscr{H}^* \otimes \mathscr{H} \to  \mathbb{C}
\] 
\[
	\left( \left\langle \phi \right| , \left| \psi \right\rangle  \right)  \mapsto \left\langle \phi \middle| \psi \right\rangle 
.\] 

In the ``wrong order'' we get the \textit{outer product}
\begin{align*}
	\mathscr{H} \otimes  \mathscr{H}^* &\longrightarrow \mathcal{B}(\mathscr{H}) \\
	\left( \left| \psi \right\rangle , \left\langle \phi \right|   \right)  &\longmapsto \left| \psi \right\rangle \left\langle \phi \right| 
.\end{align*}
\begin{align*}
	\left| \psi \right\rangle \left\langle \phi \right| : \mathscr{H} &\longrightarrow \mathscr{H} \\
	\left| x \right\rangle  &\longmapsto \left| \psi \right\rangle \left\langle \phi \middle| x \right\rangle 
.\end{align*}

In homework, we show that the outer product is an isomorphism.

\begin{example}
		On the space $\mathbb{C}^2 = \text{span}_{\mathbb{C}} \{\left| 0 \right\rangle , \left| 1 \right\rangle \} $, we have
			\[
				\left| 0 \right\rangle \left\langle 0 \right| \doteq 
					\begin{bmatrix}
						1 & 0 \\
						0 & 0 \\
					\end{bmatrix}
			.\] 
			\[
				\left| 1 \right\rangle \left\langle 0 \right| \doteq
					\begin{bmatrix}
						0 & 0 \\
						1 & 0 \\
					\end{bmatrix}
			.\] 
			Similar for $\left| 0 \right\rangle \left\langle 1 \right| $ and $\left| 1 \right\rangle \left\langle 1 \right| $.
			Thus, we have
			\[
				\begin{bmatrix}
					a & b \\
					c & d \\
				\end{bmatrix}
				= a \left| 0 \right\rangle \left\langle 0 \right| + b \left| 0 \right\rangle \left\langle 1 \right|  + c \left| 1 \right\rangle  \left\langle 0 \right|  + d \left| 1 \right\rangle \left\langle 1 \right| 
			.\] 
\end{example}

More generally, if
				\[
					A: \mathbb{C}^d \to \mathbb{C}^d
				\] 
				with $A = \left( a_{i j} \right) $, we have
				\[
					A = \sum a_{i j} \left| i \right\rangle  \left\langle j \right| 
				.\] 

\begin{example}
	$I: \mathbb{C}^d \to \mathbb{C}^d$ can be written as
	\[
		I = \sum_{i = 0}^{ d - 1} \left| i \right\rangle  \left\langle i \right| 
	.\] 
\end{example}

\begin{example}
	If $D = \operatorname{diag} \left( \lambda_0, \ldots, \lambda_{d - 1} \right) : \mathbb{C}^d \to \mathbb{C}^d$, then
	\[
		D = \sum_{i = 0}^{d - 1} \lambda_i \left| i \right\rangle  \left\langle i \right| 
	.\] 
\end{example}


\subsection{Tensor Products}

\subsubsection{Tensor product of vector spaces}


I will define tensor product $\otimes $ using basis, in which case it is also known as \textit{Kronecker product}. 

Let $V$ be a Hilbert space with orthonormal basis $\mathcal{B}_1 = \{v_1, \ldots, v_m\} $, and $W$ be another Hilbert space with $\mathcal{B}_2 = \{w_1, \ldots, w_n\} $.

Then the \textit{tensor product} is defined to be the Hilbert space $V \otimes W$ with orthonormal basis given by
\[
	\{\left| v_i \right\rangle  \otimes \left| w_j \right\rangle : 1 \le  i \le m, i \le j \le  n\} 
\] 
satisfying the following rules:
\begin{enumerate}
	\item For all $z \in \mathbb{C}$, 
		\[
			z\left( \left| v_i \right\rangle  \otimes  \left| w_j \right\rangle  \right)  = \left( z \left| v_i \right\rangle  \right) \otimes  \left| w_j \right\rangle  = \left| v_i \right\rangle  \otimes  \left( z \left| w_j \right\rangle  \right) 
		.\] 
	\item For all $1 \le i \le m$, $1 \le  j \le m$, $1 \le  k \le  n$, 
		\[
			\left( \left| v_i \right\rangle  + \left| v_j \right\rangle  \right)  \otimes  \left| w_k \right\rangle  = 
			\left| v_i \right\rangle  \otimes  \left| w_k \right\rangle  + \left| v_j \right\rangle  \otimes  \left| w_k \right\rangle 
		.\] 
		Similarly, for all $1 \le  i \le m$, $1 \le  j, k \le n$,
		\[
			\left| v_i \right\rangle  \otimes  \left( \left| w_j \right\rangle  + \left| w_k \right\rangle  \right) 
			= \left| v_i \right\rangle  \otimes  \left| w_j \right\rangle  + \left| v_k \right\rangle  \otimes  \left| w_k \right\rangle 
		.\] 
\end{enumerate}
We can check, that $\{v_i \otimes  w_k: i, k\} $ really is a basis. The inner product on $V \otimes W$ is defined so that this is a orthonormal basis.

Sometimes we might write $\left| v_i w_k \right\rangle $ for $\left| v_i \right\rangle  \otimes \left| w_k \right\rangle $.



The \textbf{curse of dimensionality}. Observe that $\operatorname{dim} V \otimes W = \left( \operatorname{dim} V \right)  \left( \operatorname{dim} W \right) .$ 
For example, for a $n$-qubit quantum memory, $\operatorname{dim} \left( \mathbb{C}^2 \right) ^{\otimes  n} = 2^n .$ This is the reason why simulating quantum systems with classical computers is in general very, very difficult.

\textbf{Binary Representation of Qubit States.}
Introduce a conventional ordering for the base states in binary order:
\[
	\left| x \right\rangle  = \sum_{i_1 = 0}^1 \ldots \sum_{i _n = 0}^1 z_{i_1, \ldots, i_n} \left| i_1 \right\rangle  \otimes  \ldots \otimes  \left| i_n \right\rangle 
	\equiv \sum z_{i_1, \ldots, i_n} \left| i_1 \ldots i_n \right\rangle 
.\] 
We can think of this as
\[
	\left| x \right\rangle = \sum_{N = 0}^{2^n - 1} z_N \left| N \right\rangle 
.\] 

\begin{example}
	Let $\left| x \right\rangle = \left| 0 \right\rangle  + \left| 1 \right\rangle  \in \mathbb{C}^2$,  $\left| y \right\rangle  = \left| 0 \right\rangle  - \left| 1 \right\rangle  \in \mathbb{C}^2$, then
	\begin{align*}
		\left| x \right\rangle \otimes \left| y \right\rangle &= \left( \left| 0 \right\rangle  + \left| 1 \right\rangle  \right) \otimes \left( \left| 0 \right\rangle  - \left| 1 \right\rangle  \right)  \\
		&= \left| 00 \right\rangle - \left| 01 \right\rangle  + \left| 10 \right\rangle  - \left| 11 \right\rangle  
	.\end{align*}
\end{example}

\subsubsection{Tensor product of operators}

For two operators $A: V_1 \to W_1, B: V_2 \to W_2$, we can define $A \otimes B: V_1 \otimes V_2 \to W_1 \otimes W_2$ such that
\[
	(A \otimes B) (v_1 \otimes v_2) = (A v_1) \otimes (B v_2)
.\] 
We can define $A \otimes B$ using a block matrix:
\[
	A \otimes B = 
	\begin{bmatrix}
		A_{1 1}B & A_{1 2} B & \ldots & A_{1 m}B \\
		A_{2 1}B & A_{2 2}B & \ldots & A_{2 m}B \\
		\vdots & \vdots & \vdots & \vdots \\
		A_{m 1}B & A_{m 2}B & \ldots & A_{m n}B \\
	\end{bmatrix}
.\] 

\begin{example}
	\[
		Z \otimes X = 
		\begin{bmatrix}
			1 & 0 \\
			0 & -1 \\
		\end{bmatrix}
		\otimes 
		\begin{bmatrix}
			0 & 1 \\
			1 & 0 \\
		\end{bmatrix}
		=
		\begin{bmatrix}
			0 & 1 & 0 & 0 \\
			1 & 0 & 0 & 0 \\
			0 & 0 & 0 & -1 \\
			0 & 0 & -1 & 0 \\
		\end{bmatrix}
		\neq X \otimes Z
	.\] 
\end{example}

\subsubsection{Simultaneous Diagonalization}



	Given $A, B: \mathscr{H} \to \mathscr{H}$, their \textit{commutator} is $[A, B] = AB - BA$. We say $A, B$ are \textit{simultaneously diagonalizable} if there exists an orthonormal basis of $\mathscr{H}$, with respect to which both $A$ and $B$ are diagonal.

\begin{theorem}[Simultaneous Diagonalization]
	\label{thm:SimultaneousDiagonalization}
	Let $A, B$ be two normal operators. $A, B$ are simultaneously diagonalizable iff they commute.
\end{theorem}
\begin{proof}
	The $\implies$ direction is trivial. Now assume that $AB = BA$. Consider eigenspace decomposition of $A$:
	\[
		A = \sum \lambda_i \left| \lambda_i \right\rangle \left\langle \lambda_i \right| 
	.\] 
	We have
	\[
		A B \left| \lambda_i \right\rangle 
	= B A \left| \lambda_i \right\rangle 
	= B \lambda_i \left| \lambda_i \right\rangle 
	= \lambda_i B \left| \lambda_i \right\rangle 
	.\] 
	Thus, $B \lambda_i$ is an eigenvector of $A$ with eigenvalue $\lambda_i$. Since $A$ has a spectral decomposition, its eigenspaces have dimension 1. Thus, $B \left| \lambda_i \right\rangle  = b_i \left| \lambda_i \right\rangle $. Combining those we have
	\[
		BA = \sum \lambda_i B \left| \lambda_i \right\rangle  \left\langle \lambda_i \right| 
			= \sum \lambda_i b_i \left| \lambda_i \right\rangle \left\langle \lambda_i \right| 
	.\] 
	We conclude that
	\[
		B = \sum b_i \left| \lambda_i \right\rangle  \left\langle \lambda_i \right| 
	.\] 
	That is, $B$ simultaneously diagonalizes with $A$.
\end{proof}


